\section{Verhalten des Modells}

Bevor die internen Mechanismen untersucht werden, lohnt sich ein einfacher Verhaltenstest: Was sagt das Modell überhaupt, wenn man es mit einem sentimentbezogenen Satz konfrontiert? Auf eine positiv formulierte Eingabe wie \texttt{The food was delicious and the service was} antwortet Pythia-410M überwiegend mit positiven Fortsetzungen wie „excellent“ oder „great“, auf eine negative Eingabe entsprechend mit „terrible“ oder „awful“. Das Modell erkennt also grundsätzlich emotionalen Kontext.

Für eine systematische Prüfung wird der gefilterte Sentiment-Challenge-Datensatz (636 Sätze, 353 negativ, 283 positiv) verwendet. An jeden Satz wird ein Prompt-Suffix angehängt, und es wird geprüft, ob unter den zehn wahrscheinlichsten nächsten Token mehr positive oder negative Lexikonwörter sind. Mit dem Suffix \texttt{ I feel very} erreicht das Modell eine Trefferquote von 75,2~Prozent, mit dem schwächeren Suffix \texttt{ Sentiment:} nur 32,1~Prozent (Tabelle~\ref{tab:cm1}). Auffällig ist dabei eine starke Neigung zur positiven Klasse: Beim zweiten Suffix sagt das Modell in 395 von 636 Fällen „positiv“ voraus und liefert in 229 weiteren Fällen gar kein erkennbares Sentiment-Token.

\begin{table}[htbp]
  \centering
  \caption{Konfusionsmatrix für den Suffix \texttt{ I feel very} (Trefferquote 75,2~\%).}
  \label{tab:cm1}
  \begin{tabular}{lrrr}
    \toprule
    & \multicolumn{3}{c}{vorhergesagt} \\
    \cmidrule(lr){2-4}
    wahr & positiv & negativ & kein Sentiment \\
    \midrule
    positiv & 259 & 24 & 0 \\
    negativ & 134 & 219 & 0 \\
    \bottomrule
  \end{tabular}
\end{table}

Tabelle~\ref{tab:catacc_short} und Abbildung~\ref{fig:catacc} zeigen die Trefferquote nach Annotationskategorie für den besseren Suffix. Eindeutige Fälle wie reine positive oder negative Sätze, Verneinungen (negated) oder Verstärkungen (amplified) werden mit über 80~Prozent gut erkannt. Am schwierigsten ist die Kategorie Sarkasmus/Ironie mit nur 46,7~Prozent.

\begin{table}[htbp]
  \centering
  \caption{Trefferquote ausgewählter Kategorien (Suffix \texttt{ I feel very}).}
  \label{tab:catacc_short}
  \begin{tabular}{lrr}
    \toprule
    Kategorie & n & Trefferquote (\%) \\
    \midrule
    positiv & 272 & 86,4 \\
    negated & 76 & 84,2 \\
    amplified & 68 & 83,8 \\
    idioms & 113 & 80,5 \\
    world-knowledge & 72 & 62,5 \\
    sarcasm/irony & 45 & 46,7 \\
    \bottomrule
  \end{tabular}
\end{table}

\begin{figure}[htbp]
  \centering
  \includegraphics[width=0.7\textwidth]{sentiment_challenge_category_accuracy_Trial1.png}
  \caption{Trefferquote je Annotationskategorie für den Suffix \texttt{ I feel very}.}
  \label{fig:catacc}
\end{figure}

Diese niedrige Quote bei Sarkasmus passt zur Erwartung, dass das Modell Ironie nicht versteht. Allerdings ist die Kategorie im Datensatz zu 93~Prozent negativ gelabelt, während das Modell ohnehin zur positiven Klasse tendiert -- ein Teil der schlechten Quote könnte also schlicht diese Schieflage widerspiegeln und nicht fehlendes Ironieverständnis. Insgesamt zeigt der Verhaltenstest: Das Modell nutzt Sentiment-Information, aber nur zuverlässig, wenn der Prompt eine sentimenttragende Fortsetzung wahrscheinlich macht. Reine Next-Token-Generierung ersetzt keine stabile Sentiment-Klassifikation -- das motiviert die folgenden Kapitel, die direkt in die internen Repräsentationen des Modells schauen.
